\documentclass[11pt]{article}

% ---------------------------
% Page layout and line spacing
% ---------------------------
\usepackage[a4paper,margin=2.5cm]{geometry}
\usepackage{setspace}
\doublespacing

% 引入行号宏包并在全局激活
\usepackage{lineno}
\linenumbers

% ---------------------------
% Encoding and fonts
% ---------------------------
\usepackage[T1]{fontenc}
\usepackage[utf8]{inputenc}
\usepackage{lmodern}
\usepackage{textcomp}
\usepackage{microtype}

% ---------------------------
% Mathematics
% ---------------------------
\usepackage{amsmath,amssymb,amsfonts,bm,mathtools}
\usepackage[document]{ragged2e}
% ---------------------------
% Figures and tables
% ---------------------------
\usepackage{graphicx}
\usepackage{booktabs}
\usepackage{array}
\usepackage{tabularx}
\usepackage{ragged2e}
\usepackage{float}
\usepackage{placeins}
\usepackage[font=small,labelfont=bf]{caption}
\usepackage{xurl}
\usepackage[hidelinks]{hyperref}

\usepackage{array}       % 支持 \newcolumntype 自定义列格式
\usepackage{longtable}   % 支持 \LTpre / \LTpost 长表格上下间距
\usepackage{ragged2e}    % 支持 \RaggedRight 智能左对齐（不产生Overfull错误）
\usepackage{tabularx}    % 支持 X 自动伸缩列


\setlength{\LTpre}{0.6em}
\setlength{\LTpost}{0.6em}
\newcolumntype{L}[1]{>{\RaggedRight\arraybackslash}p{#1}}
\newcolumntype{Y}{>{\RaggedRight\arraybackslash}X}

\renewcommand{\theequation}{S\arabic{equation}}
\renewcommand{\thetable}{S\arabic{table}}

\begin{document}

\begin{center}
{\Large \textbf{Single-cell state interpretation through neighbor evidence with Cytoneighbor}}\\[1.0em]
{\large \textbf{Supplementary Information}}
\end{center}

\vspace{1em}

\noindent\textbf{This file contains:}
\begin{itemize}
    \item Supplementary Notes S1--S8
    \item Supplementary Tables S1--S15
    \item Supplementary References
\end{itemize}

\noindent\textbf{Supplementary Notes}
\begin{itemize}
    \item Note S1. Cytoneighbor model defaults, SymRelEnc architecture and training objectives
    \item Note S2. Dynamic sampling, prototype refresh and hard-negative mining
    \item Note S3. Coarse-to-fine retrieval, distance transformation and unknown-gating annotation
    \item Note S4. Dataset sources, preprocessing, quality control and leakage control
    \item Note S5. Integration benchmark, NK--NKT boundary analysis and baseline configurations
    \item Note S6. Gene-CutMix augmentation, geometric filtering and rare-cell retrieval evaluation
    \item Note S7. DKD state-matched comparison, graph interpretation and statistical testing
    \item Note S8. Reproducibility, statistical analysis, software versions and ethics
\end{itemize}

\noindent\textbf{Supplementary Tables}
\begin{itemize}
    \item Table S1a. Dataset summary with original and quality-controlled cell and gene counts
    \item Table S1b. Dataset cell and gene counts (raw and quality-controlled)
    \item Table S2. Exact gene-processing policies and code-anchored source traces
    \item Table S3. Integration family inferential summary (paired effects)
    \item Table S4. Gene-CutMix synthesis and filtering defaults
    \item Table S5. Targeted DKD inferential comparisons
    \item Table S6. Descriptive DKD support summary
    \item Table S7. DKD targeted endpoints and descriptive inferential summary
    \item Table S8. Benchmark-specific Cytoneighbor overrides
    \item Table S9a. Software versions, exact pins and lock sources
    \item Table S9b. Baseline method hyper-parameters for annotation and integration
    \item Table S10. Reproducibility notes for software environment and repository components
    \item Table S11. Computational resources and reproducibility settings
    \item Table S12. Optional benchmark-specific marker-retention sets
    \item Table S13. LOBO fold partition definitions with reference and query cell counts
    \item Table S14. Annotation family inferential summary (random-effects meta-analysis)
    \item Table S15. Ethics, data access and code-release summary
\end{itemize}

\clearpage

\section*{Supplementary Notes}

\subsection*{Note S1. Cytoneighbor model defaults, SymRelEnc architecture and training objectives}

Unless an experiment-specific script indicated otherwise, the default SymRelEnc implementation used the following settings:
\begin{itemize}
    \item pairwise differential mode: \texttt{diff\_mode = abs}
    \item model width: \texttt{d\_model = 1024}
    \item hidden dimension: \texttt{hidden = 2048}
    \item depth: \texttt{depth = 6}
    \item dropout: \texttt{dropout = 0.1}
    \item optimizer: AdamW
    \item learning rate: \texttt{lr = 1e-3}
    \item weight decay: \texttt{weight\_decay = 1e-4}
    \item maximum epochs: \texttt{epochs = 100}
    \item warmup ratio: \texttt{warmup\_ratio = 0.05}
    \item gradient clipping: \texttt{grad\_clip = 5.0}
    \item early-stopping patience: \texttt{patience = 10}
    \item early-stopping minimum delta: \texttt{min\_delta = 1e-4}
    \item validation fraction: \texttt{val\_frac = 0.2}
    \item default monitor: \texttt{monitor = dyn\_val\_auc}
\end{itemize}

The training objective combined binary cross-entropy, rank loss and self-rank regularization. Default loss-related settings were:
\begin{itemize}
    \item \texttt{rank\_weight = 1.0}
    \item \texttt{rank\_margin = 0.5}
    \item \texttt{self\_weight = 0.2}
    \item \texttt{self\_margin = 0.5}
    \item \texttt{self\_fraction = 0.5}
    \item label smoothing disabled by default
\end{itemize}

Mixed-precision training was enabled when supported by the hardware environment, using the implementation defaults. Gradient computation used CUDA acceleration when available. These settings describe implementation defaults and should be distinguished from benchmark-specific overrides summarized in Table~S8.

\subsubsection*{S1.1 Feature alignment between query and reference}

When query and reference gene sets differ, the query matrix is aligned to the training feature list by gene symbol. The precise inference-time logic is shown below (pseudocode extracted from \texttt{cytoneighbor/model.py}, lines 219--224):

\begin{verbatim}
X_aligned = zeros(n_query_cells, n_ref_genes)
for i, gene in enumerate(ref_genes):
    if gene in query_genes:
        j = index_of(gene, query_genes)
        X_aligned[:, i] = X_query[:, j]
    else:
        X_aligned[:, i] = 0.0   # zero-filled after norm+log1p
# Query-only genes are discarded.
\end{verbatim}

\subsubsection*{S1.2 GEGLU block implementation details}

The gated residual MLP block (GEGLU-style) is defined in \texttt{cytoneighbor/model.py} (lines 437--456). Key implementation constants are:
\begin{itemize}
    \item \textbf{LayerNorm epsilon:} $10^{-5}$ (PyTorch default, not overridden).
    \item \textbf{Dropout inplace:} \texttt{False} (PyTorch default).
    \item \textbf{Residual style:} Pre-norm (LayerNorm precedes the GEGLU feed-forward sub-layer; the residual skip connects the pre-norm input to the sub-layer output).
\end{itemize}

\subsubsection*{S1.3 Validation-split construction}

The validation set (fraction 0.2) was constructed by a custom \texttt{\_group\_split} function that stratified jointly by donor and cell-type label. The split is therefore not a simple random 80/20 partition; it enforces that (i) every class is represented in both training and validation, and (ii) cells from the same donor are preferentially kept in the same split to reduce donor-level leakage. The 0.2 fraction is a global proportion, not per-class.

\subsubsection*{S1.4 Early-stopping monitor: \texttt{dyn\_val\_auc}}

Early stopping monitored \texttt{dyn\_val\_auc}, defined as the area under the ROC curve (AUC-ROC) computed on a dynamically refreshed validation pair set. Each validation pair consists of an anchor cell and a partner cell; the binary label is 1 if the two cells share the same annotated type (preferably from different donors) and 0 if they are from different types. The SymRelEnc logit for each pair is sigmoid-transformed and the AUC-ROC is computed with \texttt{sklearn.metrics.roc\_auc\_score}. The validation pair cache is refreshed every epoch with new donor-aware sampling.

\subsubsection*{S1.5 Learning-rate schedule}

The optimizer uses AdamW with an initial learning rate of $10^{-3}$. A linear warmup runs for 5 epochs ($0.05 \times 100$). After warmup, a single-cycle cosine annealing schedule decays the learning rate to $\eta_{\min}=0$ over the remaining 95 epochs. There is no cycle restart.

\subsubsection*{S1.6 Training objective}

The total loss combines binary cross-entropy, rank loss and self-rank regularization. Default loss-related settings were \texttt{rank\_weight = 1.0}, \texttt{rank\_margin = 0.5}, \texttt{self\_weight = 0.2}, \texttt{self\_margin = 0.5} and \texttt{self\_fraction = 0.5}. Label smoothing was disabled by default. Mixed-precision training was enabled when supported by the hardware environment.

\subsubsection*{S1.7 Random seeds}

The global base seed is \textbf{42}. It is propagated to Python \texttt{random}, NumPy and PyTorch (including all CUDA devices) via a \texttt{\_set\_seed} function at the start of training. Different LOBO folds do \textit{not} use different base seeds; fold-to-fold variability arises solely from the data partition. Within a single training run, epoch-specific seeds are derived as $42 + 17 \times \text{epoch}$ (training-pair sampling), $42 + 991 \times \text{epoch}$ (Gene-CutMix synthesis) and $42 + 101 \times \text{epoch} + 7$ (validation-pair refresh), preventing identical random streams across epochs while preserving reproducibility.

\subsection*{Note S2. Dynamic sampling, prototype refresh and hard-negative mining}

Training supervision was refreshed dynamically rather than fixed at initialization. The default sampler combined three components: simple label-based sampling (fraction 0.50), prototype-guided sampling (0.20) and global probe-based hard-example mining (0.30). The default number of sampled training pairs per epoch was 6000.

\subsubsection*{S2.1 Prototype computation and refresh} Prototypes were computed as the class-wise mean in the current random-projection ANN space. For each class, the $k$ samples closest to this mean were retained as prototype anchors ($k=2$ by default). Prototype pools were refreshed every 3 epochs (\texttt{proto\_update\_every = 3}).

\subsubsection*{S2.2 Global probe and hard-negative cache} Global probes were updated every epoch (\texttt{probe\_
interval = 1}) using approximate nearest-neighbor retrieval in the current feature space. For each probed anchor, 4096 candidates were retrieved and the top-32 hardest negatives were cached. Each cached entry carries a time-to-live (TTL) of 4 epochs; when TTL reaches zero the entry is purged. Hard negatives were preferentially sampled from the cache with probability 0.90; otherwise semi-hard or class-mismatched negatives were sampled by fallback rules. Rank and self-rank margins were warmed up over the early training epochs to stabilize optimization.

\subsubsection*{S2.3 Feature-space noise} Gaussian feature jitter could be injected during training with an initial standard deviation of 0.03 and subsequent epoch-wise decay. These procedures exposed the model to progressively harder positives and negatives while preserving stable optimization in sparse and heterogeneous transcriptomic settings.


\subsection*{Note S3. Coarse-to-fine retrieval, distance transformation and unknown-gating annotation}

Cytoneighbor employs a two-stage inference strategy to avoid exhaustive all-by-all pair scoring.

\subsubsection*{S3.1 Auxiliary search space and random projection}

Cells were projected into an auxiliary search space using a random projection matrix $\mathbf{P} \in \mathbb{R}^{G \times D_{\mathrm{ann}}}$. The projection matrix was initialized from a standard Gaussian distribution and orthogonalized via QR decomposition on its transpose. The auxiliary dimension was selected adaptively as
\begin{equation}
D_{\mathrm{ann}} = \min\!\left(D_{\max}, \max\!\left(D_{\min}, \mathrm{round}(rG)\right)\right),
\label{eq:aux_dim}
\end{equation}
with implementation defaults $r=0.08$, $D_{\min}=256$ and $D_{\max}=1024$. Candidate retrieval was performed using FAISS \texttt{IndexHNSWFlat} with parameters $M=32$, $\texttt{efConstruction}=200$ and $\texttt{efSearch}=128$.

\subsubsection*{S3.1.1 Worked example for auxiliary dimension selection.} For a typical human single-cell atlas with $G = 3000$ highly variable genes, the adaptive formula yields: \[ D_{\mathrm{ann}} = \min(1024, \max(256, \mathrm{round}(0.08 \times 3000))) = \min(1024, \max(256, 240)) = \min(1024, 256) = 256. \] When $G = 20{,}000$ (full genome), the formula yields $\min(1024, \max(256, 1600)) = 1024$. This adaptive sizing balances retrieval speed (lower dimension) against approximate nearest-neighbor fidelity (higher dimension captures more variance). The HNSW index parameters $M=32$ and $\texttt{efConstruction}=200$ control graph connectivity during index construction, while $\texttt{efSearch}=128$ governs the breadth of search at query time. These defaults are sufficient for reference atlases ranging from $10^{4}$ to $10^{6}$ cells.

\subsubsection*{S3.2 Coarse-candidate and refined-neighbor counts}

For graph construction through \texttt{transform()}, the default implementation retrieved up to $k_{\mathrm{coarse}}$ coarse candidates per cell before retaining 15 refined neighbors in the final graph, where
\begin{equation}
k_{\mathrm{coarse}} = \min\!\bigl(n_{\mathrm{ref}}, \max(4k_{\mathrm{final}}, 128)\bigr).
\label{eq:coarse_count}
\end{equation}
For annotation through \texttt{annotate()}, the default refined reference-neighbor count was \texttt{k\_ref = 64}, and the coarse candidate count was selected automatically as $\min(n_{\mathrm{ref}}, \max(4k_{\mathrm{ref}}, 128))$ unless overridden.

\subsubsection*{S3.3 Distance transformation for graph construction}

SymRelEnc re-scored only the retrieved candidates. For graph construction, refined logits were converted to distances using \texttt{dist\_mode = "logit\_softplus"}:
\begin{equation}
\mathrm{dist}_{ij} = \mathrm{softplus}(-\mathrm{logit}_{ij}) = \log\!\bigl(1 + \exp(-\mathrm{logit}_{ij})\bigr),
\label{eq:dist_trans}
\end{equation}
followed by numerical clamping to $[10^{-8}, 10^{6}]$. No additional distance normalization was applied. The resulting distances were written into a fuzzy neighborhood graph compatible with Scanpy downstream workflows.

\subsubsection*{S3.4 Annotation aggregation}

For label transfer, re-ranked candidate neighbors were aggregated by weighted class voting over the top-$k_{\mathrm{ref}}$ reference cells. The default weighting uses \textbf{raw logits} (\texttt{score\_mode = logit}) rather than softmax or sigmoid-transformed scores, preserving the relative magnitude of the learned pairwise evidence.

\subsubsection*{S3.5 Unknown-gating protocol in LOBO evaluation}

In the LOBO benchmark, Cytoneighbor was evaluated in label-transfer mode with \texttt{enable\_unknown = True}, \texttt{score\_mode = logit} and \texttt{temperature = 1.0}. Unknown retention used class-level and neighbor-level evidence, including top-label probability, top-versus-second-label separation, neighbor-label mixing and strongest pairwise neighbor score. This protocol should be interpreted as Cytoneighbor's inference-time unknown-gating implementation within a shared open-set-aware evaluation design, not as a fully calibrated unknown-detection benchmark for every baseline method. The open-set-aware remapping was applied at the \textbf{evaluation-script level}: any query label absent from the corresponding reference fold was remapped to \textit{Unknown} in the ground truth before metric computation, ensuring that no method was rewarded for forcing an absent label.

\subsection*{Note S4. Dataset sources, preprocessing, quality control and leakage control}

\subsubsection*{S4.1 Dataset sources} The experimental datasets utilized in this study were primarily obtained from established benchmarking and public consortium repositories. The \textbf{Human Immune} (PBMC), \textbf{Pancreas}, and \textbf{Zebrafish embryonic cells} datasets were accessed via the OpenProblems v1 repository. The \textbf{Mouse Kidney} dataset corresponds to the sci-CAR multimodal benchmark panel (GEO: GSE117089) \cite{Cao2018sciCAR}, from which only the single-cell RNA-seq modality (mod1, 11k cells) was utilized. The \textbf{Diabetic Kidney Disease} (DKD) dataset was accessed via the OpenProblems CELLxGENE census collection. The \textbf{Lung atlas} public dataset is a curated benchmarking atlas assembled by the scIB project (Theis Lab) \cite{Luecken2022IntegrationBenchmark} to evaluate data integration algorithms, comprising 32,472 cells and 15,148 genes across 16 batches and 17 cell types, downloaded from its official Figshare repository. The \textbf{Human tendon snRNA-seq} dataset was obtained directly from the CZ CELLxGENE repository (Collection ID: 579203e2-182f-47bc-8230-7aa47247e2a4, Dataset ID: 8d2e51ab-0edd-46c8-bdac-558de319d6c2), representing single-nucleus RNA sequencing of the human quadriceps tendon \cite{Tendon202x}.

Benchmark scripts consumed pre-filtered input matrices and did not re-apply a universal dataset-level \texttt{min\_counts} or \texttt{max\_mito} cell filter. Instead, preprocessing differed by benchmark driver, incorporating method-specific \texttt{filter\_genes} calls, reference-only highly variable gene (HVG) selection, and benchmark-specific split policies. Exact script-level thresholds and source traces are listed in Tables~S1a, S1b and~S2.

\subsubsection*{S4.2 Doublet removal} The released benchmark drivers did not execute an explicit doublet-removal step (e.g.\ DoubletFinder, Scrublet or SOLO). All input matrices were pre-filtered h5ad files obtained from public repositories or upstream pipelines; any doublet removal was therefore completed during the original data generation (Table~S15 lists source accessions).

\subsubsection*{S4.3 Gene-name alignment and missing-gene handling} For reference--query transfer analyses, expression features were aligned to the training feature set by gene symbol. Genes present during training but absent from a query matrix were zero-filled after normalization and log-transformation. Conversely, query-only genes not present in the trained model were discarded. This rule, enforced in \texttt{cytoneighbor/model.py} (lines 219--224), is essential for simulating realistic reference mapping in which the reference is constructed before unknown query cells are observed.

\subsubsection*{S4.4 Normalization} Unless stated otherwise, counts were normalized to a target sum of $10^{4}$ counts per cell using \texttt{sc.pp.normalize\_total} followed by \texttt{sc.pp.log1p}, yielding $\log(1 + \text{counts} \times 10^{4} / \text{total\_counts})$. No square-root, \texttt{scran} or additional scaling was applied. Cytoneighbor v7 explicitly omits a standard scaler after log-transformation.

\subsubsection*{S4.5 HVG selection and marker rescue} HVGs were computed on the reference partition only. Default HVG parameters were \texttt{flavor='seurat'}, \texttt{n\_top\_genes=3000}, with no batch stratification. In the immune-resolution stress test, the HVG step used \texttt{min\_mean=0.0125}, \texttt{max\_mean=3}, \texttt{min\_disp=0.5}, and lineage-relevant markers (e.g.\ \textit{CD3D}, \textit{KLRD1}, \textit{NCAM1}) were forcibly retained if absent from the automated HVG set (Table~S12). This marker-retention step was analysis-specific and not part of the core Cytoneighbor architecture.

\subsubsection*{S4.6 Leakage control} Preprocessing steps that could introduce leakage were defined exclusively on the reference partition. HVGs, principal components and reference dictionaries were never fitted on held-out query cells, ensuring that query data did not influence the feature space before evaluation.

\subsubsection*{S4.7 Random seeds in preprocessing} The global random seed was set to \textbf{42} for all stochastic operations, including data partitioning, UMAP, Leiden clustering and model training. Scanpy HVG selection is deterministic (statistic-based ranking) and does not require a seed.

\subsection*{Note S5. Integration benchmark, NK--NKT boundary analysis and baseline configurations}

Integration performance was evaluated using iLISI, cLISI, NMI and ARI. Because Cytoneighbor reconstructs a pairwise-score-based neighborhood graph rather than directly optimizing a coordinate-space latent representation, integration behaviour was interpreted using these complementary summaries rather than through any single metric.

\subsubsection*{S5.1 LISI computation}

iLISI and cLISI were computed with the \texttt{scib} package (version listed in Table~S9a) using \texttt{perplexity = 30}. NMI and ARI were computed against the annotated cell-type labels after Leiden clustering on the respective latent or graph representations.

\subsubsection*{S5.2 Baseline hyper-parameters}

Table~S9b lists the exact hyper-parameters used for each baseline method. Key settings include:
\begin{itemize}
    \item \textbf{Seurat:} \texttt{FindTransferAnchors} with \texttt{npcs = 30}, \texttt{dims = 1:30}, \texttt{k.anchor = 5}, \texttt{k.filter = 200}; \texttt{TransferData} with \texttt{k.weight = 50}.
    \item \textbf{SingleR:} \texttt{de.method = "wilcox"}, \texttt{fine.tune = TRUE}, \texttt{prune = TRUE}.
    \item \textbf{CellTypist:} Reference-retrained model (not a pre-trained atlas model), \texttt{max\_iter = 100}, \texttt{feature\_selection = True}, majority voting disabled.
    \item \textbf{scVI:} \texttt{n\_layers = 2}, \texttt{n\_latent = 30}, \texttt{dropout\_rate = 0.1}, \texttt{gene\_likelihood = "zinb"}, trained for up to 400 epochs with early-stopping patience 45.
    \item \textbf{Harmony:} \texttt{theta = 2.0}, \texttt{lambda = 1.0}, \texttt{sigma = 0.1}, \texttt{max\_iter\_harmony = 10}.
    \item \textbf{BBKNN:} \texttt{n\_pcs = 50}, \texttt{neighbors\_within\_batch = 3}.
\end{itemize}

\subsubsection*{S5.3 NK--NKT boundary stress test}

To examine increasingly subtle subtype boundaries, a difficulty gradient was defined in the Human Immune dataset spanning CD4$^{+}$ T cells, CD8$^{+}$ T cells, NK cells and NKT cells. In this analysis, lineage-relevant markers were retained after feature selection to avoid discarding weakly detected but biologically informative subtype markers (Table~S12). Performance was assessed using global and class-specific F1 scores, confusion matrices, targeted NKT-to-NK misclassification rates and focused local-neighborhood visualizations in the NK--NKT region.

This boundary was chosen because it is both continuous and separable: close enough to be vulnerable to over-smoothing, yet distinct enough that spillover into NK or CD8 has different biological meanings. In the residual-error summary, the denominator is fixed to all true NKT cells in each method, and errors are decomposed into NK, CD8 and Other destinations so that local boundary blur can be distinguished from adjacent-lineage confusion and broader non-specific spillover. The NK route reflects artificial continuity across a boundary that should remain readable, whereas the CD8 route reflects a more distant lineage confusion with different downstream biological consequences. Because the method-specific zooms are shown in each representation's native UMAP coordinate space, the supporting comparison is based on local neighbor structure, marker anchoring and residual error routes rather than on absolute axis ranges. This supporting analysis should therefore be interpreted as a local immune-boundary readability study rather than as a generic benchmark comparison.

\subsubsection*{S5.4 Integration metric interpretation across datasets}

Integration performance was evaluated across the Human Immune, Pancreas and Diabetic Kidney Disease datasets. In the Human Immune dataset, Harmony achieved the highest iLISI (3.616), whereas Cytoneighbor showed the strongest NMI (0.857) and ARI (0.786). In Pancreas, Harmony again achieved the highest iLISI (2.646), while Cytoneighbor retained the strongest NMI (0.849) and ARI (0.725). In DKD, Harmony maximized iLISI (3.875), whereas Cytoneighbor showed the strongest NMI (0.868) and ARI (0.871). These summaries support the interpretation that stronger mixing alone does not automatically provide the most biologically readable representation, particularly once local immune structure and downstream disease interpretation are considered.


\subsection*{Note S6. Gene-CutMix augmentation, geometric filtering and rare-cell retrieval evaluation}

Gene-CutMix was applied only during training and was formalized as a constrained synthesis operator over compatible parent cells.

\subsubsection*{S6.1 Synthesis operator}

In the default two-parent mode, for compatible parents $(\mathbf{x}_i, \mathbf{x}_j)$, a contiguous mask $\mathbf{M} \in \{0,1\}^G$ was sampled and
\begin{equation}
\tilde{\mathbf{x}} = \mathbf{M} \odot \mathbf{x}_i + (\mathbf{1}-\mathbf{M}) \odot \mathbf{x}_j,\quad r\sim \mathcal{U}(0.35,0.65),
\label{eq:cutmix_two}
\end{equation}
where $r$ controls the exchanged gene fraction. In the optional three-parent mode (default probability \texttt{synth\_p\_three\_parent}=0.25), three contiguous segments with default proportions $(0.3,0.4,0.3)$ are concatenated according to
\begin{equation}
\tilde{\mathbf{x}} = \mathbf{M}_1 \odot \mathbf{x}_{p_1} + \mathbf{M}_2 \odot \mathbf{x}_{p_2} + \mathbf{M}_3 \odot \mathbf{x}_{p_3},\quad \mathbf{M}_1+\mathbf{M}_2+\mathbf{M}_3=\mathbf{1}.
\label{eq:cutmix_three}
\end{equation}
Compatibility is enforced before synthesis by a rule set implemented in code: class-label consistency is required; donor-aware and lineage-aware constraints are additionally applied when those annotations are available.

\subsubsection*{S6.2 Geometric filter and NaN fallback}

Accepted synthetic cells also passed a geometric filter. For each donor--class key $(d,c)$, the implementation computed a centre $\boldsymbol{\mu}_{d,c}$, intra-class radius $\rho_{d,c}$ (95th percentile distance to centre), and nearest other-class centre distance $\delta_{d,c}$ within donor. A synthetic cell was retained only if
\begin{equation}
\|\tilde{\mathbf{x}} - \boldsymbol{\mu}_{d,c}\|_2 \le \min(\alpha\rho_{d,c},\, \gamma\delta_{d,c}),
\label{eq:geom_filter}
\end{equation}
with defaults $\alpha = \gamma = 1.0$. If radius or margin statistics were undefined (e.g.\ insufficient support), the corresponding bound was treated as $+\infty$ in the implementation. If \textit{both} thresholds were $+\infty$, the synthetic cell was accepted (no geometric constraint). The exact implementation resides in \texttt{cytoneighbor/model.py} (lines 615--631).

\subsubsection*{S6.3 Rare-cell mechanism analysis and interpretation}

For the rare-cell mechanism analysis, pancreatic $\delta$ cells were aggressively downsampled in the reference ($\texttt{downsample\_fraction}=0.01$, donor-stratified random sampling), while a held-out set of real $\delta$ cells was reserved for retrieval testing. Cross-partition Recall@$K$ was computed by concatenating reference and query cells, constructing neighbors on the combined set, and retaining only neighbors belonging to the real reference partition for evaluation. A safety buffer of $N_{\mathrm{query}}$ (the number of query cells) was added to the requested neighbor count so that, after removal of query self-hits, each query cell still retained at least $K$ genuine reference neighbors.

Under this design, the main interpretation remains that held-out rare query cells either do or do not encounter real rare reference neighbors near the top of the ranking. Improvements in the selected run therefore reflect stronger ranked neighbor evidence rather than direct augmentation of the test-time retrieval pool by synthetic cells. In the selected run used for the main figure, the first correct rare reference neighbor appeared earlier, rank-resolved neighborhood composition shifted toward more $\delta$-cell evidence within a fixed early-neighborhood window, and Recall@$K$ remained positive across the evaluated range. This is read not only as retrieval improvement, but as earlier recovery of $\delta$-cell neighbor evidence before standard annotation visibly fails.

For the \textit{SST}/\textit{INS} marker panel, each held-out $\delta$ query contributed the marker value of its retrieved real reference neighbor at a fixed rank $k$, and mean $\pm$ SE was then computed across queries. Query-level aggregation was chosen so that each rare query contributed equally to the marker summary and query-specific variation from more densely sampled regions could not dominate the estimate. The corresponding biological interpretation is that recovered neighbors become more \textit{SST}-high and less \textit{INS}-high, indicating restoration of $\delta$-cell neighbor evidence rather than only earlier retrieval. In other words, the earliest failure under extreme scarcity is not simply loss of final label accuracy, but loss of the neighbor evidence needed to read $\delta$ cells as $\delta$ cells.

Gene-CutMix should be interpreted narrowly as a training-time stabilizer for severely underrepresented local neighborhoods. Synthetic profiles were not inserted into the retrieval database, and the observed benefit should not be taken as evidence for biological intermediate states or a uniquely correct rare-cell graph structure.

\subsubsection*{S6.4 Synthesis volume}

The number of synthetic cells generated per epoch was dynamic. For each flagged difficult type, the target multiplicity was $\texttt{synth\_mult\_per\_dt}=1.2$ times the real count, capped at $\texttt{synth\_cap\_per\_dt}=600$ and over-generated by a factor of 1.5 before geometric filtering. Synthetic cells typically accounted for 5--15\% of the training pairs sampled in a given epoch.

\subsubsection*{S6.5 Synthesis volume statistics.} In the pancreatic rare-cell analysis with $\delta$-cell downsampling ($\texttt{downsample\_fraction}=0.01$), the remaining reference count was: \[ n_{\delta,\mathrm{ref}} \approx \text{\texttt{[estimated from Pancreas rare-cell subset]}}. \] Gene-CutMix generated $\approx 1.2 \times n_{\delta,\mathrm{ref}} \times 1.5$ candidates before geometric filtering. These synthetic profiles were used only during training-pair sampling and did not appear in the inference-time retrieval database. Across all cell types in the combined Pancreas training set, synthetic cells contributed approximately 5--15\% of the total training pairs sampled per epoch. In the Human Immune dataset used for annotation benchmarks, synthesis volume ranged from approximately 5\% to 15\% depending on the LOBO fold, reflecting variable rare-cell counts across batches. These statistics can be extracted from training logs available in the repository.

\subsection*{Note S7. DKD state-matched comparison, graph interpretation and statistical testing}

The relational model was trained using all healthy reference cells rather than proximal-tubule cells alone, preserving multi-lineage negative pairs during training. After training, healthy and diseased proximal-tubule cells were concatenated into a joint cell-by-gene expression matrix and processed using one Cytoneighbor graph-construction step.

\subsubsection*{S7.1 Clustering and root selection}

Leiden clustering was used to define coarse PT states at \textbf{resolution = 0.5}. This value was selected after an exploratory search over $\{0.3, 0.4, 0.5, 0.6, 0.8, 1.0\}$ that favoured a three-cluster solution with biologically coherent marker profiles and high modularity. The root state for downstream graph interpretation was selected without marker-based pre-assignment: the cluster containing the largest absolute number of healthy proximal-tubule cells was used. This reduced circularity because zonation identities and marker-based names were assigned only after graph construction.

\subsubsection*{S7.2 PAGA and diffusion pseudotime}

PAGA connectivities were computed with $\texttt{n\_neighbors = 15}$ and Euclidean distance on the Cytoneighbor embedding. Weak connections below a threshold of 0.01 were filtered for visualization. Diffusion pseudotime was used only as an approximate ordering of transcriptional change within the graph; neither PAGA nor DPT was treated as proof of an \textit{in vivo} lineage trajectory.

\subsubsection*{S7.3 Gene Ontology enrichment}

GO enrichment was performed with \texttt{gseapy} (Enrichr API) against the 2021 GO gene sets (Biological Process, Molecular Function and Cellular Component). The background gene set consisted of all expressed genes (detection rate $>1\%$ of cells) rather than the full genome. Significance thresholds were nominal $p < 0.05$ and Benjamini--Hochberg adjusted $p < 0.05$.

\subsubsection*{S7.4 Functional and injury scores}

The composite functional score (based on \textit{SLC5A2} and \textit{LRP2}) and injury score (based on \textit{VCAM1} and \textit{HAVCR1}) were computed by Z-score standardizing each marker's $\log(1+x)$ expression across cells, then averaging the Z-scores within each cell:
\begin{equation}
\text{Score}_i = \frac{1}{|M|} \sum_{g \in M} z_{ig}, \qquad z_{ig} = \frac{x_{ig} - \bar{x}_g}{s_g}.
\label{eq:score_formula}
\end{equation}
This ensures equal weighting across markers regardless of absolute expression scale.

\subsubsection*{S7.5 Within-state statistical testing}

Within-state comparisons of SLC5A2 expression between healthy and DKD cells were tested with one-sided Mann--Whitney $U$ tests using \texttt{scipy.stats.mannwhitneyu} with \texttt{alternative = 'greater'} (testing DKD $>$ Healthy). The family of comparisons comprised the three proximal-tubule states (S1/S2-like, S3-like, Injured), and $p$-values were adjusted within this family using the Benjamini--Hochberg procedure at $\alpha = 0.05$ (Table~S5). Exact comparison definitions and descriptive support are summarized in Tables~S5 and~S6.

\subsubsection*{S7.6 Additional caveats on DKD state interpretation}

Healthy-enriched proximal-tubule states were interpreted only after graph construction and cluster definition. Marker-based naming was therefore performed \textit{post hoc} rather than used to define the graph itself. In particular, \textit{SLC5A1} was used as an independent indicator of S3-like identity within the functional proximal-tubule states, whereas \textit{SLC5A2}, \textit{LRP2} and related epithelial transport genes were interpreted in combination with graph topology and cell-of-origin composition. Elevated \textit{SLC5A1} signal in injured proximal-tubule cells was interpreted descriptively as retention of segment-associated gene-expression features and may reflect relative enrichment of injured cells that still retain S3-linked transcriptional features, rather than serving as proof of lineage origin or segment-specific susceptibility.

PAGA connectivity was used to assess whether the three-state organization was more consistent with convergence onto a shared injured state than with a simple linear arrangement. Diffusion pseudotime was then used only as an approximate ordering within this joint graph. The DKD analysis should therefore be interpreted as graph-supported transcriptional organization within the analyzed cohort rather than as a definitive account of \textit{in vivo} lineage progression. Preserved neighborhood structure is treated as the condition under which segment-specific organization and injury-associated remodeling can be interpreted together, not as proof of a single true disease trajectory.

\subsection*{Note S8. Reproducibility, statistical analysis, software versions and ethics}

\subsubsection*{S8.1 Statistical analysis registry and annotation metric reporting}

A complete registry of all formal statistical tests, including test name, contrast groups, sample sizes, test statistics, $p$-values, effect sizes and correction methods, is provided in Table~S13. Macro-F1 was used as the primary annotation metric because it weights classes equally and is therefore less dominated by abundant populations than overall accuracy. For LOBO annotation, scores were computed after open-set-aware ground-truth remapping of query labels absent from the corresponding reference fold. For integration analyses, iLISI, cLISI, NMI and ARI were treated as complementary summaries of batch mixing and biological conservation.

For the Cytoneighbor LOBO benchmark, the reported results correspond to fold-specific models trained with Gene-CutMix enabled and evaluated in label-transfer mode with \texttt{enable\_unknown = True}, \texttt{score\_mode = logit} and \texttt{temperature = 1.0}. Query labels absent from the corresponding reference fold were remapped in the ground truth to \textit{Unknown} before metric computation. Accordingly, the reported annotation metrics should be interpreted as robustness to neighborhood truncation under incomplete reference coverage within a shared evaluation protocol rather than as a claim of equivalent calibrated unknown-detection capability across all compared methods. The full fold-wise annotation results are available in the repository results directory.

\subsubsection*{S8.2 Inferential summary methodology}

For inferential reporting, the primary estimand was a paired effect size $\Delta = \text{metric}_{\text{Cytoneighbor}} - \text{metric}_{\text{baseline}}$, computed at fold level for annotation and at dataset level for integration. For each pre-specified family (\textbf{Annotation}, \textbf{Integration}, \textbf{DKD targeted}), we report effect size, 95\% CI, nominal $p$, and Benjamini--Hochberg FDR-adjusted $q$ within family. Annotation and integration used paired two-sided $t$-based summaries on the exported fold/dataset units; DKD targeted inference retained the pre-specified one-sided Mann--Whitney $U$ formulation for the within-state \textit{SLC5A2} directional test with alternative hypothesis DKD $>$ healthy and nominal $\alpha=0.05$.

\subsubsection*{S8.3 Random seeds and stochastic settings}

All stochastic stages used a base seed of \textbf{42}. The specific seed propagation is documented in Table~S11. Because \texttt{deterministic = False} in the PyTorch backend (performance priority), exact bitwise reproducibility across different GPU architectures or CUDA versions is not guaranteed; however, statistical conclusions remain stable across runs.

\subsubsection*{S8.4 Computational resources}

Training was performed on NVIDIA GPUs with CUDA support across high-performance computing environments, including local workstations, university servers, and virtual machines. Peak GPU memory usage was governed by an adaptive batch-size calculator with a default budget of 2048~MB, which scaled dynamically depending on atlas size. Computational runtime and resource allocation (VRAM, CPU cores, system RAM) varied across experiments based on the specific hardware deployed and the scale of the corresponding dataset fold.

\subsubsection*{S8.5 Software environment and code freeze}

Package versions are pinned in \texttt{repro/locks/conda-list-smn.txt}, \texttt{repro/locks/pip-freeze-smn
.txt} and \texttt{repro/locks/environment-smn.yml}. The Git freeze tag \texttt{submission-v1-freeze} marks the exact code state used for all results.

\subsubsection*{S8.6 Ethics and data availability}

All human datasets were obtained from public repositories (GEO, ArrayExpress, CELLxGENE, Figshare) and comprise de-identified, secondary-use data. Original studies obtained institutional review board approval and informed consent as cited in the primary publications. No new human samples were collected for this study. Processed data, model checkpoints and result tables are available under the accession and repository details listed in Table~S15. Patient-level identifiable metadata (exact age, sex, clinical identifiers) were excluded from the release.


\clearpage

\section*{Supplementary Tables}

\noindent\textbf{Supplementary Data guide.} All raw data files, fold-wise annotation results, integration benchmark summaries, statistical assets, model checkpoints and result tables are available from the GitHub repository and from the corresponding author upon reasonable request. Machine-readable outputs include:
\begin{itemize}
    \item Fold-wise LOBO annotation metrics (accuracy, macro-F1, micro-F1, weighted-F1, NMI, AMI, ARI)
    \item Integration benchmark summaries (iLISI, cLISI, NMI, ARI) across Human Immune, Pancreas and Diabetic Kidney Disease
    \item Statistical inference assets (\texttt{repro/statistics/annotation\_paired\_effects.csv}, \texttt{integr-
    ation\_paired\_effects.csv}, \texttt{dkd\_targeted\_effects.csv})
    \item QC threshold manifest (\texttt{repro/qc\_threshold\_manifest.csv})
    \item Rare-cell retrieval assets (\texttt{results/mechanism\_analysis/recall\_curve\_long.csv}, \texttt{rare\_q-
    uery\_summary.csv})
    \item DKD characterization summaries (\texttt{results/disease\_discovery/experiment\_results.json})
\end{itemize}

\FloatBarrier

\begin{table}[htbp]
\centering
\renewcommand{\thetable}{S2}
\caption{\textbf{Exact gene-processing policies and code-anchored source traces.}}
\label{tab:s2_preprocessing_anchor}
\scriptsize
\begin{tabularx}{\textwidth}{@{}L{0.22\textwidth} Y L{0.28\textwidth}@{}}
\toprule
\textbf{Dataset} & \textbf{Gene/HVG preprocessing (exact)} & \textbf{Source anchor} \\
\midrule
Human Immune (integration + resolution) & Integration: \texttt{filter\_genes(min\_cells=10)} + \texttt{normalize\_total(1e4)} + \texttt{log1p} + \texttt{HVG(3000)}; Resolution: \texttt{HVG(min\_mean=0.0125,max\_mean=3,min\_disp=0.5)} with mandatory-marker rescue & \texttt{experiments/Rebatch/a-main\_remove.py}; \texttt{experiments/benchmark\_resolution.py} \\
Pancreas (rare-cell) & \texttt{normalize\_total(1e4)} + \texttt{log1p}; reference-only \texttt{HVG(0.0125,3,0.5)}; rare downsampling \texttt{downsample\_fraction=0.01}, \texttt{n\_keep}=max(2,$\dots$) & \texttt{experiments/mechanism\_analysis.py} \\
Diabetic Kidney Disease (joint graph) & \texttt{min\_cells}=max(10,0.001$\times n_{\mathrm{obs}}$) and \texttt{gene\_ct}$\ge$10; \texttt{normalize\_total(1e4)} + \texttt{log1p} + \texttt{HVG(3000)} & \texttt{experiments/disease\_discovery.py} \\
Lung atlas public (LOBO) & Method-specific wrappers: \texttt{filter\_genes(min\_cells=3/10/20)} and reference-only HVG (\texttt{2000} or \texttt{3000}) as required & Method-specific wrappers in \texttt{experiments/Annotation/} and \texttt{experiments/Rebatch/} \\
Mouse kidney (mod1, LOBO) & Cytoneighbor wrapper: reference \texttt{normalize\_total(1e4)} + \texttt{log1p} + \texttt{HVG(3000)}; CellTypist wrapper: \texttt{filter\_genes(min\_cells=10)} & \texttt{experiments/Annotation/Cytoneighbor.py}; \texttt{experiments/Annotation/} baselines \\
Zebrafish embryonic cells (LOBO) & scANVI/scArches wrappers: \texttt{filter\_genes(min\_cells=10/20)}, \texttt{HVG(3000)}; SciBet wrapper: \texttt{HVG(2000)} & \texttt{experiments/Annotation/} and \texttt{experiments/Rebatch/} baseline wrappers \\
Human tendon snRNA-seq & Same LOBO wrapper policy as above (method-specific \texttt{filter\_genes} and reference-only HVG where required) & \texttt{repro/qc\_threshold\_manifest.csv} \\
\bottomrule
\end{tabularx}
\end{table}

\FloatBarrier

\begin{table}[htbp]
\centering
\renewcommand{\thetable}{S1b}
\caption{\textbf{Dataset cell and gene counts (raw and quality-controlled).}}
\label{tab:s1b_dataset_counts}
\scriptsize
\begin{tabularx}{\textwidth}{@{}L{0.20\textwidth} L{0.12\textwidth} L{0.12\textwidth} L{0.12\textwidth} Y@{}}
\toprule
\textbf{Dataset} & \textbf{Cells (raw)} & \textbf{Cells (QC)} & \textbf{Genes (raw)} & \textbf{QC filtering description} \\
\midrule
Human Immune & 33,506 & 33,506 & 12,303 & Pre-filtered upstream; varies by batch \\
Pancreas & 16,382 & 16,382 & 18,771 & Multi-study compendium; pre-filtered \\
Diabetic Kidney Disease & 39,176 & 39,176 & 27,980 & GSE131882; dynamic gene threshold \\
Lung atlas public & 32,472 & 32,472 & 15,148 & scIB benchmark atlas (16 batches, 17 cell types); runtime load \\
Mouse Kidney (mod1) & 11,296 & 11,296 & 29,013 & Benchmark panel (sci-CAR); mod1 RNA only \\
Zebrafish embryonic cells & 26,022 & 26,022 & 25,258 & Embryo atlas panel; runtime load \\
Human tendon snRNA-seq & 12,808 & 12,808 & 61,455 & CELLxGENE; sample/donor stratified \\
\bottomrule
\end{tabularx}
\vspace{0.5em}

{\raggedright\scriptsize\textit{Notes.} Cell and gene counts extracted from input h5ad files using \texttt{extract\_dataset\_statistics.py}. Raw counts reflect the input matrix shape; QC counts reflect cells retained after method-specific filtering. For compendium datasets assembled from multiple studies, counts represent the combined total after merging. Actual statistics derived from script output.\par}
\end{table}

\FloatBarrier

\begin{table}[htbp]
\centering
\renewcommand{\thetable}{S2}
\caption{\textbf{Exact gene-processing policies and code-anchored source traces.}}
\label{tab:s2_preprocessing_anchor}
\scriptsize
\begin{tabularx}{\textwidth}{@{}L{0.22\textwidth} Y L{0.28\textwidth}@{}}
\toprule
\textbf{Dataset} & \textbf{Gene/HVG preprocessing (exact)} & \textbf{Source anchor} \\
\midrule
Human Immune (integration + resolution) & Integration: \texttt{filter\_genes(min\_cells=10)} + \texttt{normalize\_total(1e4)} + \texttt{log1p} + \texttt{HVG(3000)}; Resolution: \texttt{HVG(min\_mean=0.0125,max\_mean=3,min\_disp=0.5)} with mandatory-marker rescue & \texttt{Rebatch/a-main\_remove.py:2732--2735}; \texttt{benchmark\_resolution.py:197--216} \\
Pancreas (rare-cell) & \texttt{normalize\_total(1e4)} + \texttt{log1p}; reference-only \texttt{HVG(0.0125,3,0.5)}; rare downsampling \texttt{downsample\_fraction=0.01}, \texttt{n\_keep}=max(2,$\dots$) & \texttt{mechanism\_analysis.py:63--66,104--105, 113,127--128} \\
Diabetic Kidney Disease (joint graph) & \texttt{min\_cells}=max(10,0.001$\times n_{\mathrm{obs}}$) and \texttt{gene\_ct}$\ge$10; \texttt{normalize\_total(1e4)} + \texttt{log1p} + \texttt{HVG(3000)} & \texttt{disease\_discovery.py:273--274, 289--291,312--315} \\
Lung atlas public (LOBO) & Method-specific wrappers: \texttt{filter\_genes(min\_cells=3/10/20)} and reference-only HVG (\texttt{2000} or \texttt{3000}) as required & \texttt{Annotation/a-main\_seurat\_cv-ls.py:72}; \texttt{a-main\_scanvi\_cv-ls.py:208,251}; \texttt{a-main\_scarches-scanvi-joint.py:266} \\
Mouse kidney (mod1, LOBO) & Cytoneighbor wrapper: reference \texttt{normalize\_total(1e4)} + \texttt{log1p} + \texttt{HVG(3000)}; CellTypist wrapper: \texttt{filter\_genes(min\_cells=10)} & \texttt{a-main\_Cytoneighbor\_cv-ls.py:69--70}; \texttt{a-main\_celltypist\_cv-ls.py:244} \\
Zebrafish embryonic cells (LOBO) & scANVI/scArches wrappers: \texttt{filter\_genes(min\_cells=10/20)}, \texttt{HVG(3000)}; SciBet wrapper: \texttt{HVG(2000)} & \texttt{a-main\_scanvi\_cv-ls.py:154,208,251}; \texttt{a-main\_scarches\_cv-ls.py:24,44,63}; \texttt{a-main\_scibet\_cv-ls.py:14,30,44} \\
Human tendon snRNA-seq & Same LOBO wrapper policy as above (method-specific \texttt{filter\_genes} and reference-only HVG where required) & \texttt{repro/qc\_threshold\_manifest.csv} \\
\bottomrule
\end{tabularx}
\end{table}

\FloatBarrier

\begin{table}[htbp]
\centering
\renewcommand{\thetable}{S3}
\caption{\textbf{Integration family inferential summary.}}
\label{tab:s3_integration_effects}
\scriptsize
\begin{tabularx}{\textwidth}{@{}L{0.10\textwidth} L{0.15\textwidth} c L{0.12\textwidth} L{0.20\textwidth} L{0.10\textwidth} Y@{}}
\toprule
\textbf{Metric} & \textbf{Baseline} & \textbf{$n$} & \textbf{Effect $\Delta$} & \textbf{95\% CI} & \textbf{$p$} & \textbf{BH-$q$} \\
\midrule
iLISI & Unintegrated & 3 & $+0.719$ & $[0.436,\ 1.001]$ & 0.0083 & 0.0857 \\
iLISI & Harmony & 3 & $-1.083$ & $[-2.160,\ -0.006]$ & 0.0495 & 0.1282 \\
iLISI & scVI & 3 & $-0.526$ & $[-1.605,\ 0.553]$ & 0.1710 & 0.2280 \\
iLISI & BBKNN & 3 & $-0.109$ & $[-0.673,\ 0.456]$ & 0.4952 & 0.5282 \\
cLISI & Unintegrated & 3 & $+0.004$ & $[-0.040,\ 0.049]$ & 0.7188 & 0.7188 \\
cLISI & Harmony & 3 & $-0.014$ & $[-0.056,\ 0.027]$ & 0.2813 & 0.3462 \\
cLISI & scVI & 3 & $-0.036$ & $[-0.082,\ 0.009]$ & 0.0738 & 0.1476 \\
cLISI & BBKNN & 3 & $-0.006$ & $[-0.034,\ 0.022]$ & 0.4265 & 0.4875 \\
NMI & Unintegrated & 3 & $+0.178$ & $[0.098,\ 0.258]$ & 0.0107 & 0.0857 \\
NMI & Harmony & 3 & $+0.088$ & $[-0.006,\ 0.181]$ & 0.0561 & 0.1282 \\
NMI & scVI & 3 & $+0.094$ & $[0.034,\ 0.154]$ & 0.0215 & 0.0859 \\
NMI & BBKNN & 3 & $+0.096$ & $[0.024,\ 0.167]$ & 0.0290 & 0.0929 \\
ARI & Unintegrated & 3 & $+0.451$ & $[0.187,\ 0.715]$ & 0.0180 & 0.0859 \\
ARI & Harmony & 3 & $+0.233$ & $[-0.181,\ 0.648]$ & 0.1365 & 0.1985 \\
ARI & scVI & 3 & $+0.209$ & $[-0.141,\ 0.559]$ & 0.1239 & 0.1983 \\
ARI & BBKNN & 3 & $+0.257$ & $[-0.088,\ 0.602]$ & 0.0853 & 0.1516 \\
\bottomrule
\end{tabularx}
\vspace{0.5em}

{\raggedright\scriptsize\textit{Notes.} Effect is defined as Cytoneighbor $-$ baseline. BH-FDR correction was applied within the Integration family. Full machine-readable outputs are available in the repository statistics directory.\par}
\end{table}

\FloatBarrier

\begin{table}[htbp]
\centering
\renewcommand{\thetable}{S4}
\caption{\textbf{Gene-CutMix synthesis and filtering defaults used in the released implementation.}}
\label{tab:s4_genecutmix_defaults}
\scriptsize
\begin{tabularx}{\textwidth}{@{}L{0.22\textwidth} L{0.15\textwidth} Y L{0.22\textwidth}@{}}
\toprule
\textbf{Parameter / rule} & \textbf{Default} & \textbf{Formal meaning} & \textbf{Source anchor} \\
\midrule
\texttt{synth\_enable} & \texttt{True} & Enable train-time synthetic generation & \texttt{cytoneighbor/model.py:941} \\
\texttt{synth\_refresh\_every} & 1 & Rebuild synthesis pool each epoch interval & \texttt{cytoneighbor/model.py:942,1146} \\
\texttt{synth\_mult\_per\_dt} / \texttt{synth\_cap\_per\_dt} & 1.2 / 600 & Per difficult-type target multiplicity and cap & \texttt{cytoneighbor/model.py:943--944,1149} \\
\texttt{synth\_overgen\_factor} & 1.5 & Over-generate candidates before filtering & \texttt{cytoneighbor/model.py:945,1150} \\
\texttt{synth\_ratio\_range} & (0.35, 0.65) & Two-parent exchanged gene fraction interval & \texttt{cytoneighbor/model.py:946,531--541} \\
\texttt{synth\_p\_three\_parent} & 0.25 & Probability of three-parent synthesis mode & \texttt{cytoneighbor/model.py:947,545--566} \\
\texttt{synth\_parts\_three} & (0.3, 0.4, 0.3) & Three contiguous segment proportions & \texttt{cytoneighbor/model.py:948,545--566} \\
\texttt{synth\_alpha}, \texttt{synth\_gamma} & 1.0, 1.0 & Radius and margin scaling in acceptance inequality & \texttt{cytoneighbor/model.py:949--950} \\
Intra-class radius $\rho_{d,c}$ & 95th percentile distance & $\rho_{d,c}=Q_{0.95}\left(\|\mathbf{x}-\boldsymbol{\mu}_{d,c}\|_2\right)$, fallback $+\infty$ for low support & \texttt{cytoneighbor/model.py:571--594} \\
Nearest other-class margin $\delta_{d,c}$ & min center distance & $\delta_{d,c}=\min_{c'\neq c}\|\boldsymbol{\mu}_{d,c}-\boldsymbol{\mu}_{d,c'}\|_2$, fallback $+\infty$ when undefined & \texttt{cytoneighbor/model.py:571--608} \\
Acceptance filter & inequality rule & Keep synthetic sample iff $\|\tilde{\mathbf{x}}-\boldsymbol{\mu}_{d,c}\|_2 \le \min(\alpha \rho_{d,c},\gamma \delta_{d,c})$ & \texttt{cytoneighbor/model.py:615--631} \\
\bottomrule
\end{tabularx}
\end{table}

\FloatBarrier

\begin{table}[htbp]
\centering
\renewcommand{\thetable}{S5}
\caption{\textbf{Targeted DKD inferential comparisons.}}
\label{tab:s5_dkd_stats}
\scriptsize
\begin{tabularx}{\textwidth}{@{}L{0.20\textwidth} L{0.20\textwidth} L{0.15\textwidth} L{0.20\textwidth} Y@{}}
\toprule
\textbf{Comparison} & \textbf{Cell subset} & \textbf{Direction tested} & \textbf{Statistical test} & \textbf{Adjustment} \\
\midrule
\textit{SLC5A2} Healthy vs DKD & S1/S2-like PT & DKD $>$ Healthy & one-sided Mann--Whitney $U$ & BH-FDR within 3-cluster family \\
\textit{SLC5A2} Healthy vs DKD & S3-like PT & DKD $>$ Healthy & one-sided Mann--Whitney $U$ & BH-FDR within 3-cluster family \\
\textit{SLC5A2} Healthy vs DKD & Injured PT & DKD $>$ Healthy & one-sided Mann--Whitney $U$ & BH-FDR within 3-cluster family \\
\bottomrule
\end{tabularx}
\end{table}

\begin{table}[htbp]
\centering
\renewcommand{\thetable}{S6}
\caption{\textbf{Descriptive DKD support summary.}}
\label{tab:s6_dkd_descriptive}
\scriptsize
\begin{tabularx}{\textwidth}{@{}L{0.24\textwidth} L{0.28\textwidth} Y@{}}
\toprule
\textbf{Comparison} & \textbf{Status / source} & \textbf{Interpretive note} \\
\midrule
Functional program score (SLC5A2/LRP2) by state & Cluster-level state summaries (0 vs 2); descriptive state-level summary from \texttt{experiment\_results.json} & Functional PT (S1/S2-like) $>$ Injured PT (1.962 vs 1.222); descriptive support; inferential $p$ was not retained in exported summary file \\
Injury program score (VCAM1/HAVCR1) by state & Cluster-level state summaries & Injured PT $>$ Functional PT (0.705 vs 0.079); descriptive support \\
\bottomrule
\end{tabularx}
\end{table}

\FloatBarrier

\begin{table}[htbp]
\centering
\renewcommand{\thetable}{S7}
\caption{\textbf{DKD targeted endpoints and descriptive inferential summary.}}
\label{tab:s7_dkd_targeted_registry}
\scriptsize
\begin{tabularx}{\textwidth}{@{}L{0.12\textwidth} L{0.25\textwidth} Y L{0.12\textwidth} L{0.20\textwidth}@{}}
\toprule
\textbf{Main-text location} & \textbf{Endpoint} & \textbf{Definition / test} & \textbf{Effect or lock} & \textbf{Status / note} \\
\midrule
Fig.~5h & \texttt{SLC5A2\_S1S2 (DKD>Healthy)} & One-sided Mann--Whitney $U$ (pre-specified directional) & nominal $p$ reported in main text & Targeted test specified in pipeline; multiplicity-adjusted $q$ and confidence interval not retained in exported \texttt{experiment\_results.json} \\
Fig.~5g & \texttt{FuncScore (C0-C2)} & Descriptive state-level contrast & $+0.732$ & Derived from exported DKD characterization summary \\
Fig.~5g & \texttt{InjuryScore (C2-C0)} & Descriptive state-level contrast & $+0.623$ & Derived from exported DKD characterization summary \\
\bottomrule
\end{tabularx}
\vspace{0.5em}

{\raggedright\scriptsize\textit{Notes.} Table S7 summarizes currently retained DKD targeted outputs available from released artifacts. Machine-readable files are available in the repository statistics directory. For the within-state \textit{SLC5A2} directional comparison, nominal $p$ values are reported in the main text; multiplicity-adjusted $q$ values and confidence intervals were not retained in the exported JSON artifact.\par}
\end{table}

\FloatBarrier

\begin{table}[htbp]
\centering
\renewcommand{\thetable}{S8}
\caption{\textbf{Benchmark-specific Cytoneighbor overrides relative to implementation defaults.}}
\label{tab:s8_overrides}
\scriptsize
\begin{tabularx}{\textwidth}{@{}L{0.20\textwidth} L{0.22\textwidth} L{0.22\textwidth} Y@{}}
\toprule
\textbf{Benchmark} & \textbf{Key override} & \textbf{Value} & \textbf{Rationale} \\
\midrule
LOBO annotation & \texttt{enable\_unknown} & \texttt{True} & Cytoneighbor inference used unknown gating within the shared open-set-aware evaluation protocol \\
LOBO annotation & \texttt{score\_mode} & \texttt{logit} & label-transfer evaluation \\
LOBO annotation & \texttt{temperature} & 1.0 & default score scaling \\
Integration / graph definition & final neighbors & 15 & graph construction for downstream Scanpy analyses \\
Rare-cell retrieval analysis & neighbor request buffer & analysis-specific & ensure at least $K$ evaluable reference neighbors \\
DKD joint graph & training population & all healthy cells & preserve multi-lineage negatives during training \\
\bottomrule
\end{tabularx}
\end{table}

\begin{table}[htbp]
\centering
\renewcommand{\thetable}{S9a}
\caption{\textbf{Software versions, exact pins and lock sources.}}
\label{tab:s9a_software_versions}
\scriptsize
\begin{tabularx}{\textwidth}{@{}L{0.25\textwidth} L{0.25\textwidth} Y@{}}
\toprule
\textbf{Component} & \textbf{Exact pin} & \textbf{Lock / source} \\
\midrule
Python runtime & 3.10.18 & \texttt{repro/locks/conda-list-smn.txt} \\
PyTorch & 2.8.0+cu129 & \texttt{repro/locks/conda-list-smn.txt} \\
Scanpy / AnnData / FAISS & 1.10.3 / 0.11.4 / 1.8.0 & \texttt{repro/locks/conda-list-smn.txt} \\
scvi-tools / scArches / CellTypist & 1.3.3 / 0.6.1 / 1.6.3 & \texttt{repro/locks/conda-list-smn.txt} \\
R / Seurat / SingleR & 4.4.1 / 5.3.0 / 2.8.0 & \texttt{repro/locks/r-session-smn.txt} \\
SciBet baseline & wrapper-script pin & \texttt{experiments/Annotation/NetWork/run\_scibet.R} \\
gseapy (GO enrichment) & version in conda lock & \texttt{repro/locks/conda-list-smn.txt} \\
scib (integration metrics) & version in conda lock & \texttt{repro/locks/conda-list-smn.txt} \\
Git freeze tag & \texttt{submission-v1-freeze} & git tag in repository root \\
Package locks & explicit conda + pip + env export & \texttt{repro/locks/*.txt,yml} \\
Container recipes & Docker + Apptainer definitions & \texttt{repro/container/*} \\
\bottomrule
\end{tabularx}
\vspace{0.5em}

{\raggedright\scriptsize\textit{Notes.} For baseline method-specific hyper-parameters, see Table~S9b.\par}
\end{table}

\FloatBarrier

\begin{table}[htbp]
\centering
\renewcommand{\thetable}{S9b}
\caption{\textbf{Baseline method hyper-parameters for annotation and integration.}}
\label{tab:s9b_baseline_params}
\scriptsize
\begin{tabularx}{\textwidth}{@{}L{0.20\textwidth} L{0.28\textwidth} Y@{}}
\toprule
\textbf{Method} & \textbf{Parameter} & \textbf{Value / Configuration} \\
\midrule
\multicolumn{3}{@{}l@{}}{\textit{scVI (scvi-tools)}} \\
& \texttt{n\_layers} & 2 \\
& \texttt{n\_latent} & 30 \\
& \texttt{dropout\_rate} & 0.1 \\
& \texttt{gene\_likelihood} & "zinb" (Zero-Inflated Negative Binomial) \\
& \texttt{max\_epochs} & 400 \\
& early stopping patience & 45 \\
& training data & Raw counts restored from \texttt{layers['counts']} before model setup \\
\midrule
\multicolumn{3}{@{}l@{}}{\textit{Seurat TransferData}} \\
& \texttt{NormalizeData} & \texttt{verbose = FALSE} \\
& \texttt{FindVariableFeatures} & \texttt{selection.method = "vst"}, \texttt{nfeatures = 2000} \\
& \texttt{ScaleData} & Applied to variable features \\
& \texttt{RunPCA} & Applied to variable features, \texttt{verbose = FALSE} \\
& \texttt{FindTransferAnchors.dims} & 1:30 \\
& \texttt{FindTransferAnchors.reduction} & "rpca" (reciprocal PCA) \\
& \texttt{FindTransferAnchors.features} & Variable features (2000 genes) \\
& \texttt{TransferData.dims} & 1:30 \\
& execution mode & Single-threaded (sequential) due to memory constraints \\
\midrule
\multicolumn{3}{@{}l@{}}{\textit{SingleR}} \\
& \texttt{trainSingleR} & Default parameters \\
& \texttt{classifySingleR} & Default parameters \\
& pruning strategy & \texttt{pruned.labels} used (ambiguous cells pruned to NA) \\
& input data & log-normalized counts (\texttt{logcounts} assay) \\
\midrule
\multicolumn{3}{@{}l@{}}{\textit{CellTypist}} \\
& training mode & Reference-retrained (not pre-trained atlas model) \\
& \texttt{check\_expression} & \texttt{True} (validates input is log1p-transformed) \\
& \texttt{n\_jobs} & -1 (use all available CPU cores) \\
& \texttt{feature\_selection} & \texttt{True} if $n_{\text{genes}} \ge 300$, else \texttt{False} \\
& prediction mode & \texttt{mode = 'best match'} (default, recommended) \\
& majority voting & Disabled (\texttt{majority\_voting = False}) \\
& data preprocessing & NaN labels removed from reference; NaN values in X replaced with 0.0 \\
\midrule
\multicolumn{3}{@{}l@{}}{\textit{SciBet}} \\
& core function & \texttt{SciBet\_R()} from \texttt{scibetR} package \\
& input format & Dense data frame (genes $\times$ cells) \\
& implementation & R wrapper script: \texttt{experiments/Annotation/NetWork/run\_scibet.R} \\
\midrule
\multicolumn{3}{@{}l@{}}{\textit{Harmony}} \\
& \texttt{theta} & 2.0 (diversity clustering penalty) \\
& \texttt{lambda} & 1.0 (ridge regression penalty) \\
& \texttt{sigma} & 0.1 (width of soft k-means clusters) \\
& \texttt{max\_iter\_harmony} & 10 (maximum correction iterations) \\
\midrule
\multicolumn{3}{@{}l@{}}{\textit{BBKNN (Batch-Balanced K-Nearest neighbors)}} \\
& \texttt{n\_pcs} & 50 \\
& \texttt{neighbors\_within\_batch} & 3 \\
\bottomrule
\end{tabularx}
\vspace{0.5em}

{\raggedright\scriptsize\textit{Notes.} All baseline methods were evaluated under a shared open-set-aware protocol: query labels absent from the reference fold were remapped to \textit{Unknown} in the ground truth before metric computation. Code anchors: LOBO baseline scripts are located in \texttt{experiments/Annotation/} (e.g., Seurat, SingleR, CellTypist, SciBet, and \texttt{NetWork/} R wrappers); Integration baselines (scVI, Harmony, BBKNN) are executed via \texttt{experiments/benchmark\_resolution.py}.\par}
\end{table}

\begin{table}[htbp]
\centering
\renewcommand{\thetable}{S10}
\caption{\textbf{Reproducibility notes for software environment and repository components.}}
\label{tab:s10_software_notes}
\scriptsize
\begin{tabularx}{\textwidth}{@{}L{0.24\textwidth} Y@{}}
\toprule
\textbf{Component} & \textbf{Notes} \\
\midrule
Python runtime & Primary runtime in conda env \texttt{smn} \\
PyTorch & Core Cytoneighbor training/inference dependency \\
Scanpy / AnnData / FAISS & Graph analysis, data structure and HNSW backend \\
scvi-tools / scArches / CellTypist & Baseline stack in the same pinned environment \\
R / Seurat / SingleR & R baseline versions pinned from the execution environment \\
SciBet baseline & Package-level SciBet version is not separately pinned; wrapper is pinned by repository freeze tag \\
gseapy & Python wrapper for Enrichr GO API; used in DKD analysis \\
scib & Integration-metric suite (iLISI, cLISI, NMI, ARI) \\
Git freeze tag & Added in this revision to replace prior ``no-tag'' state \\
Package locks & SHA256 entries recorded in \texttt{repro/repro\_manifest.csv} \\
Container recipes & Recipe files are hash-pinned; image digests are obtained at build time as documented in \texttt{repro/container/README.md} \\
\bottomrule
\end{tabularx}
\end{table}

\FloatBarrier

\begin{table}[htbp]
\centering
\renewcommand{\thetable}{S11}
\caption{\textbf{Computational resources and reproducibility settings.}}
\label{tab:s11_compute_resources}
\scriptsize
\begin{tabularx}{\textwidth}{@{}L{0.24\textwidth} Y@{}}
\toprule
\textbf{Resource / Setting} & \textbf{Specification} \\
\midrule
GPU model & NVIDIA GPU with CUDA support \\
GPU memory & Variable depending on deployment environment (batch size adapted automatically) \\
CUDA version & 11.x--12.x (compatible with PyTorch 2.x) \\
Peak GPU memory usage & Dynamically scaled per LOBO fold (governed by adaptive batch-size calculator with default budget 2048 MB) \\
Training time per LOBO fold & Dependent on dataset scale and hardware (early stopping typically triggered after 30--50 epochs) \\
CPU cores & Variable (sufficient for data loading and preprocessing) \\
System RAM & Variable (scaled to dataset size) \\
Operating system & Linux (Ubuntu 20.04+ or CentOS 7+ recommended) \\
Base random seed & 42 (propagated to Python \texttt{random}, NumPy, PyTorch including CUDA) \\
Epoch-specific seeds & Training pairs: $42 + 17 \times \text{epoch}$; Gene-CutMix: $42 + 991 \times \text{epoch}$; Validation pairs: $42 + 101 \times \text{epoch} + 7$ \\
PyTorch deterministic mode & \texttt{deterministic = False} (performance priority; bitwise reproducibility not guaranteed across GPU architectures) \\
Mixed precision training & Enabled when hardware-supported (default PyTorch AMP implementation) \\
Parallelization & Training: single-GPU per fold; Inference: batch-parallelized candidate retrieval via FAISS \\
Disk space requirement & $\approx$5--20 GB per dataset (raw h5ad + preprocessed + model checkpoints + result tables) \\
\bottomrule
\end{tabularx}
\vspace{0.5em}

{\raggedright\scriptsize\textit{Notes.} Ranges reflect hardware variability across different execution environments. Exact package versions are listed in Table~S10.\par}
\end{table}

\FloatBarrier

\begin{table}[htbp]
\centering
\renewcommand{\thetable}{S12}
\caption{\textbf{Optional benchmark-specific marker-retention sets.}}
\label{tab:s12_marker_sets}
\scriptsize
\begin{tabularx}{\textwidth}{@{}L{0.22\textwidth} Y@{}}
\toprule
\textbf{Analysis} & \textbf{Example retained markers} \\
\midrule
Immune resolution boundary analysis & \textit{CD4}, \textit{CD8A}, \textit{CD8B}, \textit{CD3D}, \textit{CD3E}, \textit{CD3G}, \textit{NCAM1}, \textit{KLRD1}, \textit{FCGR3A}, \textit{IL7R}, \textit{CCR7}, \textit{S100A4}, \textit{FOXP3}, \textit{GZMB}, \textit{PRF1} (retained when absent from HVGs in the fine-boundary analysis pipeline) \\
DKD zonation validation & \textit{SLC5A1} (used to distinguish the S3-like functional state from the S1/S2-like state), \textit{SLC5A2}, \textit{LRP2} (functional PT program), \textit{VCAM1}, \textit{HAVCR1} (injury program) for post hoc state interpretation \\
\bottomrule
\end{tabularx}
\end{table}

\begin{table}[htbp]
\centering
\renewcommand{\thetable}{S13}
\caption{\textbf{LOBO fold partition definitions with reference and query cell counts.}}
\label{tab:s13_lobo_partitions}
\scriptsize
\begin{tabularx}{\textwidth}{@{}L{0.22\textwidth} L{0.22\textwidth} L{0.12\textwidth} L{0.12\textwidth} Y@{}}
\toprule
\textbf{Dataset} & \textbf{Held-out batch} & \textbf{Ref. cells} & \textbf{Query cells} & \textbf{Notes} \\
\midrule
Human Immune & Villani & varies & varies & Myeloid/lymphoid stress-test fold \\
Human Immune & Oetjen\_A & varies & varies & Public benchmark fold \\
Human Immune & Oetjen\_P & varies & varies & Public benchmark fold \\
Human Immune & Oetjen\_U & varies & varies & Public benchmark fold \\
Human Immune & Freytag & varies & varies & Public benchmark fold \\
Human Immune & 10X & varies & varies & Public benchmark fold \\
Human Immune & Sun\_sample1\_CS & varies & varies & Public benchmark fold \\
Pancreas & Smarter & varies & varies & Endocrine stress-test fold \\
Pancreas & Baron 2016 & varies & varies & Public benchmark fold \\
Pancreas & Segerstolpe 2016 & varies & varies & Public benchmark fold \\
Pancreas & Xin 2016 & varies & varies & Public benchmark fold \\
Pancreas & Lawlor 2017 & varies & varies & Public benchmark fold \\
Lung atlas public & multiple & varies & varies & Per-fold counts vary \\
Mouse Kidney (mod1) & multiple & varies & varies & Per-fold counts vary \\
Zebrafish embryo & multiple & varies & varies & Per-fold counts vary \\
\bottomrule
\end{tabularx}
\vspace{0.5em}

{\raggedright\scriptsize\textit{Notes.} Reference (Ref.) cells were used for model training and retrieval database construction; Query cells were held out for evaluation. Complete fold-wise annotation metrics (accuracy, macro-F1, micro-F1, weighted-F1, NMI, AMI, ARI) are available in the repository results directory.\par}
\end{table}

\FloatBarrier

\begin{table}[htbp]
\centering
\renewcommand{\thetable}{S14}
\caption{\textbf{Annotation family inferential summary (Macro-F1 random-effects meta-analysis across datasets).}}
\label{tab:s14_annotation_meta}
\scriptsize
\begin{tabularx}{\textwidth}{@{}L{0.15\textwidth} c L{0.12\textwidth} L{0.20\textwidth} L{0.10\textwidth} L{0.10\textwidth} Y@{}}
\toprule
\textbf{Baseline} & \textbf{$k$ datasets} & \textbf{Effect $\Delta$} & \textbf{95\% CI} & \textbf{$p$} & \textbf{BH-$q$} & \textbf{$\tau^2$} \\
\midrule
CellTypist & 6 & $-0.011$ & $[-0.070,\ 0.048]$ & 0.7129 & 0.7129 & 0.0048 \\
SciBet & 7 & $+0.052$ & $[-0.099,\ 0.202]$ & 0.5025 & 0.7129 & 0.0404 \\
Seurat & 7 & $+0.026$ & $[-0.024,\ 0.077]$ & 0.3073 & 0.6145 & 0.0036 \\
SingleR & 7 & $+0.063$ & $[0.003,\ 0.122]$ & 0.0402 & 0.1206 & 0.0060 \\
scANVI & 6 & $-0.087$ & $[-0.153,\ -0.021]$ & 0.0101 & 0.0603 & 0.0055 \\
scArches & 6 & $-0.007$ & $[-0.041,\ 0.027]$ & 0.6802 & 0.7129 & 0.0009 \\
\bottomrule
\end{tabularx}
\vspace{0.5em}

{\raggedright\scriptsize\textit{Notes.} Effect is defined as Cytoneighbor $-$ baseline on Macro-F1. BH-FDR correction was applied within the Annotation family. Full dataset-level and fold-level results are available in the repository statistics directory.\par}
\end{table}

\FloatBarrier

\begin{table}[htbp]
\centering
\renewcommand{\thetable}{S15}
\caption{\textbf{Ethics, data access and code-release summary.}}
\label{tab:s15_ethics_data}
\scriptsize
\begin{tabularx}{\textwidth}{@{}L{0.22\textwidth} L{0.38\textwidth} Y@{}}
\toprule
\textbf{Item} & \textbf{Details} & \textbf{Access / identifier} \\
\midrule
Human Immune (PBMC) & De-identified public compendium & OpenProblems v1: immune\_cells \\
Pancreas & De-identified public islet atlases & OpenProblems v1: pancreas \\
DKD & De-identified public kidney data & OpenProblems CELLxGENE census: dkd \\
Lung atlas & De-identified public lung atlas & scIB Figshare repository \cite{Luecken2022IntegrationBenchmark} \\
Mouse kidney & De-identified public benchmark & OpenProblems v1 multimodal: scicar\_mouse\_kidney \cite{Cao2018sciCAR} \\
Zebrafish embryo & De-identified public atlas & OpenProblems v1: zebrafish \\
Human tendon & De-identified snRNA-seq & CELLxGENE: 8d2e51ab-0edd-46c8-bdac-558de319d6c2 \cite{Tendon202x} \\
Ethics waiver & Secondary analysis of de-identified data & No new human subjects; exempt from additional IRB review \\
Processed data & Preprocessed h5ad, checkpoints, result tables & Available from the corresponding author upon reasonable request \\
Code repository & Full source, locks, end-to-end scripts & Available from the corresponding author upon reasonable request \\
Patient metadata & Excluded from release & Only anonymous donor IDs and cell types retained \\
\bottomrule
\end{tabularx}
\end{table}

\FloatBarrier

\clearpage

\renewcommand{\refname}{Supplementary References}
\begin{thebibliography}{99}

\bibitem{Regev2017HCA}
Regev, A. et al. The Human Cell Atlas. \textit{eLife} \textbf{6}, e27041 (2017).

\bibitem{Luecken2019BestPractices}
Luecken, M. D. \& Theis, F. J. Current best practices in single-cell RNA-seq analysis: a tutorial. \textit{Mol. Syst. Biol.} \textbf{15}, e8746 (2019).

\bibitem{Lahnemann2020Challenges}
L\"ahnemann, D. et al. Eleven grand challenges in single-cell data science. \textit{Genome Biol.} \textbf{21}, 31 (2020).

\bibitem{Kiselev2019Challenges}
Kiselev, V. Y., Andrews, T. S. \& Hemberg, M. Challenges in unsupervised clustering of single-cell RNA-seq data. \textit{Nat. Rev. Genet.} \textbf{20}, 273--282 (2019).

\bibitem{Heumos2023BestPractices}
Heumos, L. et al. Best practices for single-cell analysis across modalities. \textit{Nat. Rev. Genet.} \textbf{24}, 550--572 (2023).

\bibitem{Aran2019SingleR}
Aran, D. et al. Reference-based analysis of lung single-cell sequencing reveals a transitional profibrotic macrophage. \textit{Nat. Immunol.} \textbf{20}, 163--172 (2019).

\bibitem{Stuart2019Seurat}
Stuart, T. et al. Comprehensive integration of single-cell data. \textit{Cell} \textbf{177}, 1888--1902 (2019).

\bibitem{DominguezConde2022CellTypist}
Dom\'{\i}nguez Conde, C. et al. Cross-tissue immune cell analysis reveals tissue-specific features in humans. \textit{Science} \textbf{376}, eabl5197 (2022).

\bibitem{Abdelaal2019ComparisonAnnotation}
Abdelaal, T. et al. A comparison of automatic cell identification methods for single-cell RNA sequencing data. \textit{Genome Biol.} \textbf{20}, 194 (2019).

\bibitem{Pasquini2021AnnotationReview}
Pasquini, G., Rojo Arias, J. E., Sch\"afer, P. \& Busskamp, V. Automated methods for cell type annotation on scRNA-seq data. \textit{Comput. Struct. Biotechnol. J.} \textbf{19}, 961--969 (2021).

\bibitem{Haghverdi2018BatchEffects}
Haghverdi, L., Lun, A. T. L., Morgan, M. D. \& Marioni, J. C. Batch effects in single-cell RNA-sequencing data are corrected by matching mutual nearest neighbors. \textit{Nat. Biotechnol.} \textbf{36}, 421--427 (2018).

\bibitem{Luecken2022IntegrationBenchmark}
Luecken, M. D. et al. Benchmarking atlas-level data integration in single-cell genomics. \textit{Nat. Methods} \textbf{19}, 41--50 (2022).

\bibitem{Lake2023KidneyAtlas}
Lake, B. B. et al. An atlas of healthy and injured cell states and niches in the human kidney. \textit{Nature} \textbf{619}, 585--594 (2023).

\bibitem{Korsunsky2019Harmony}
Korsunsky, I. et al. Fast, sensitive and accurate integration of single-cell data with Harmony. \textit{Nat. Methods} \textbf{16}, 1289--1296 (2019).

\bibitem{Lopez2018scVI}
Lopez, R., Regier, J., Cole, M. B., Jordan, M. I. \& Yosef, N. Deep generative modeling for single-cell transcriptomics. \textit{Nat. Methods} \textbf{15}, 1053--1058 (2018).

\bibitem{Polanski2020BBKNN}
Pola\'nski, K. et al. BBKNN: fast batch alignment of single cell transcriptomes. \textit{Bioinformatics} \textbf{36}, 964--965 (2020).

\bibitem{HubertArabie1985ARI}
Hubert, L. \& Arabie, P. Comparing partitions. \textit{J. Classif.} \textbf{2}, 193--218 (1985).

\bibitem{Vinh2010NMI}
Vinh, N. X., Epps, J. \& Bailey, J. Information theoretic measures for clusterings comparison: variants, properties, normalization and correction for chance. \textit{J. Mach. Learn. Res.} \textbf{11}, 2837--2854 (2010).

\bibitem{Villani2017Immune}
Villani, A.-C. et al. Single-cell RNA-seq reveals new types of human blood dendritic cells, monocytes, and progenitors. \textit{Science} \textbf{356}, eaah4573 (2017).

\bibitem{Muraro2016Pancreas}
Muraro, M. J. et al. A single-cell transcriptome atlas of the human pancreas. \textit{Cell Syst.} \textbf{3}, 385--394 (2016).

\bibitem{Baron2016Pancreas}
Baron, M. et al. A single-cell transcriptomic map of the human and mouse pancreas reveals inter- and intra-cell population structure. \textit{Cell Syst.} \textbf{3}, 346--360 (2016).

\bibitem{Segerstolpe2016Pancreas}
Segerstolpe, \AA. et al. Single-cell transcriptome profiling of human pancreatic islets in health and type 2 diabetes. \textit{Cell Metab.} \textbf{24}, 593--607 (2016).

\bibitem{Xin2016Pancreas}
Xin, Y. et al. RNA sequencing of single human islet cells reveals type 2 diabetes genes. \textit{Cell Metab.} \textbf{24}, 608--615 (2016).

\bibitem{Lawlor2017Pancreas}
Lawlor, N. et al. Single-cell transcriptomes identify human islet cell signatures and reveal cell-type-specific expression changes in type 2 diabetes. \textit{Genome Res.} \textbf{27}, 208--222 (2017).

\bibitem{Muto2021Kidney}
Muto, Y. et al. Single cell transcriptional and chromatin accessibility profiling redefine cellular heterogeneity in the adult human kidney. \textit{Nat. Commun.} \textbf{12}, 2190 (2021).

\bibitem{Wilson2022DKD}
Wilson, P. C. et al. Multimodal single cell sequencing implicates chromatin accessibility and genetic background in diabetic kidney disease progression. \textit{Nat. Commun.} \textbf{13}, 5253 (2022).

\bibitem{Gerhardt2021PTStates}
Gerhardt, L. M. S. et al. Single-nuclear transcriptomics reveals diversity of proximal tubule cell states in a dynamic response to acute kidney injury. \textit{Proc. Natl. Acad. Sci. USA} \textbf{118}, e2026684118 (2021).

\bibitem{Johnson2019FAISS}
Johnson, J., Douze, M. \& J\'egou, H. Billion-scale similarity search with GPUs. \textit{IEEE Trans. Big Data} \textbf{7}, 535--547 (2019).

\bibitem{Malkov2018HNSW}
Malkov, Y. A. \& Yashunin, D. A. Efficient and robust approximate nearest neighbor search using hierarchical navigable small world graphs. \textit{IEEE Trans. Pattern Anal. Mach. Intell.} \textbf{42}, 824--836 (2020).

\bibitem{Wolf2018Scanpy}
Wolf, F. A., Angerer, P. \& Theis, F. J. SCANPY: large-scale single-cell gene expression data analysis. \textit{Genome Biol.} \textbf{19}, 15 (2018).

\bibitem{Xu2021scANVI}
Xu, C. et al. Probabilistic harmonization and annotation of single-cell transcriptomics data with deep generative models. \textit{Mol. Syst. Biol.} \textbf{17}, e9620 (2021).

\bibitem{Lotfollahi2022scArches}
Lotfollahi, M. et al. Mapping single-cell data to reference atlases by transfer learning. \textit{Nat. Biotechnol.} \textbf{40}, 121--130 (2022).

\bibitem{Traag2019Leiden}
Traag, V. A., Waltman, L. \& van Eck, N. J. From Louvain to Leiden: guaranteeing well-connected communities. \textit{Sci. Rep.} \textbf{9}, 5233 (2019).

\bibitem{Wolf2019PAGA}
Wolf, F. A. et al. PAGA: graph abstraction reconciles clustering with trajectory inference through a topology preserving map of single cells. \textit{Genome Biol.} \textbf{20}, 59 (2019).

\bibitem{Haghverdi2016DPT}
Haghverdi, L., B\"uttner, M., Wolf, F. A., Buettner, F. \& Theis, F. J. Diffusion pseudotime robustly reconstructs lineage branching. \textit{Nat. Methods} \textbf{13}, 845--848 (2016).

\bibitem{Cao2018sciCAR}
Cao, J. et al. Joint profiling of chromatin accessibility and gene expression in thousands of single cells. \textit{Science} \textbf{361}, 1380--1385 (2018).

\bibitem{Tendon202x}
Mimpen, J. Y. et al. Exploring cellular changes in ruptured human quadriceps tendons at single-cell resolution. \textit{J. Physiol.} \textbf{603}, 4542--4560 (2025).

\end{thebibliography}

\end{document}